\documentclass[runningheads]{llncs}

\usepackage[T1]{fontenc}
\usepackage{graphicx}
\usepackage{listings}
\usepackage{xcolor}
\usepackage{pdfpages}
\usepackage{hyperref}

\lstset{
    language=HTML,
    basicstyle=\ttfamily\small,
    keywordstyle=\color{blue},
    stringstyle=\color{orange},
    commentstyle=\color{gray},
    breaklines=true,
    frame=single,
    numbers=left,
    numberstyle=\tiny\color{gray},
    showstringspaces=false,
}

\begin{document}

\title{Módulo 1: Los Pookies}

\author{Allamand Francisco \and
Mateo Notti Echave \and
Facundo Dip \and
Pedro Elizalde}



\institute{Universidad Nacional de Cuyo\\
\email{franallamand05@gmail.com}}

\maketitle

\begin{abstract}
En este trabajo se presentan los conocimientos adquiridos durante el Módulo 1, enfocados en el uso de lenguajes de marcado y herramientas digitales orientadas al desarrollo académico y profesional. Se abordan conceptos fundamentales de HTML, Markdown y LaTeX, junto con su aplicación práctica en la estructuración de documentos y contenidos web. Además, se incorporan recursos de referencia y la integración de documentación personal como parte del aprendizaje.
\end{abstract}

\section{Introducción}
A lo largo de este módulo se produjo una transición desde el uso básico de herramientas digitales hacia la adopción de metodologías más estructuradas y profesionales. El aprendizaje de lenguajes de marcado como HTML, Markdown y LaTeX permitió comprender cómo organizar información de manera clara, reutilizable y estandarizada \cite{github, markdown, latex}.

Este informe tiene como objetivo mostrar la aplicación concreta de estos lenguajes mediante ejemplos prácticos, destacando su utilidad tanto en el desarrollo web como en la producción de documentos técnicos. Asimismo, se evidencia la importancia del trabajo colaborativo y el uso de herramientas como GitHub en entornos académicos vinculados a la ingeniería \cite{git}.

\section{Uso avanzado de HTML}

\subsection{Ejemplo completo}
\begin{lstlisting}[language=HTML, caption=Ejemplo completo de HTML]
<!DOCTYPE html>
<html lang="es">
<head>
    <meta charset="UTF-8">
    <title>Ejemplo Completo</title>
</head>
<body>

<h1>Bienvenido</h1>
<p><b>Texto en negrita</b> y <i>texto en italica</i></p>

<hr>

<h2>Lista</h2>
<ul>
<li>Elemento 1</li>
<li>Elemento 2</li>
</ul>

<h2>Tabla</h2>
<table border="1">
<tr>
<th>Nombre</th>
<th>Edad</th>
</tr>
<tr>
<td>Juan</td>
<td>20</td>
</tr>
</table>

<h2>Enlaces</h2>
<a href="https://www.google.com">Ir a Google</a>

<h2>Imagen</h2>
<img src="imagen.jpg" alt="Ejemplo">

</body>
</html>
\end{lstlisting}

\subsection{Segundo ejemplo HTML}
\begin{lstlisting}[language=HTML, caption=Segundo ejemplo HTML]
<!DOCTYPE html>
<html>
<body>

<h1>Pagina simple</h1>

<p>Ir a una seccion interna:</p>
<a href="#seccion">Seccion</a>

<h2 id="seccion">Seccion interna</h2>
<p>Contenido de la seccion.</p>

</body>
</html>
\end{lstlisting}

\section{Lenguajes de Marcado}

\subsection{Markdown}
\begin{lstlisting}[caption=Ejemplo Markdown]
# Titulo
## Subtitulo

Texto en **negrita** y *italica*

- Item 1
- Item 2

[Link](https://github.com)
\end{lstlisting}

\subsection{LaTeX}
\begin{lstlisting}[language=TeX, caption=Ejemplo LaTeX]
\section{Seccion}
Texto con cita \cite{latex}.

\begin{itemize}
\item Punto 1
\item Punto 2
\end{itemize}
\end{lstlisting}

\section{Cheatsheets}

Para la correcta implementación de los lenguajes de marcado y el uso de las herramientas del módulo, se consultaron las siguientes guías rápidas y documentación oficial:

\begin{itemize}
\item \href{https://www.markdownguide.org/cheat-sheet/}{Markdown Cheat Sheet} \cite{markdown}
\item \href{https://www.w3schools.com}{HTML Reference} \cite{html}
\item \href{https://www.overleaf.com/learn}{LaTeX Guide} \cite{latex}
\item \href{https://education.github.com}{GitHub Guide} \cite{git}
\end{itemize}

\section{Conclusión}
La integración de herramientas como HTML, Markdown y LaTeX durante este primer módulo no solo representa la adquisición de habilidades técnicas, sino que establece una base metodológica fundamental para la comunicación en ingeniería. La distinción entre el contenido y su representación visual permite un manejo mucho más eficiente de la información, facilitando la creación de documentación técnica que sea, al mismo tiempo, precisa y estandarizada.

La experiencia práctica con estos lenguajes, sumada a la gestión de entornos colaborativos en plataformas como GitHub, refleja las exigencias del ámbito profesional contemporáneo. En este contexto, la capacidad de estructurar informes técnicos complejos y de mantener un control de versiones riguroso se traduce en una ventaja competitiva directa. En última instancia, el dominio de estas tecnologías de marcado asegura que el flujo de trabajo sea escalable y adaptable a las diversas necesidades de la industria y la investigación académica.

\section{Anexos}

% PDFs de CV
\includepdf[pages=-, pagecommand={\thispagestyle{plain}}]{facundodipcv.pdf}

\bibliographystyle{splncs04}
\bibliography{bibliografia}

\end{document}